\fichatitulo{24 · RECUPERAÇÃO E GERAÇÃO}{EACL · 2021}{Fusion-in-Decoder}
{\small Izacard; Grave. \textit{Leveraging Passage Retrieval with Generative Models for Open Domain Question Answering}. EACL · 2021. \href{https://aclanthology.org/2021.eacl-main.74.pdf}{Fonte consultada}.\par}

\campo{Problema e objetivo}
Como combinar evidência distribuída entre passagens recuperadas?

\campo{Principal contribuição · síntese interpretativa}
Propõe processar passagens separadamente no codificador e fundir suas representações no decodificador para gerar respostas.

\campo{Método / natureza da evidência}
Experimentos de perguntas de domínio aberto, incluindo Natural Questions e TriviaQA, variando a quantidade de passagens.

\campo{Resultado ou argumento principal}
A arquitetura consegue explorar até cem passagens nas condições estudadas, ampliando a combinação de evidências.

\campo{Excertos-chave · seleção literal}
\begin{itemize}
\excerto{1}{combine evidence from multiple passages}{Resumo.}
\excerto{2}{up to one hundred passages}{Introdução.}
\end{itemize}

\campo{Limitações e análise crítica}
Análise da ficha: ganhos dependem da arquitetura, treinamento e tarefa. O resultado não autoriza fixar cem trechos como configuração ótima de qualquer sistema RAG.

\campo{Aproveitamento na dissertação · proposta}
Fundamentação: distinguir fusão no decodificador de simples concatenação em um prompt; testar quantidade de contexto.

\registro{Fonte consultada nesta preparação: Resumo e apresentação do método. Chave: \texttt{izacardGrave2021passageRetrieval}.}
